%%%%%%%%%%%%%%%%%%%%%%%%%%%%%%%%%%%%%%%%%
% Modelo Latex para pôster do SBrT'24
%
% v1: Diego da Silva de Medeiros - IFSC
% v2: Leonardo Lira Ramalho - UFPA
%%%%%%%%%%%%%%%%%%%%%%%%%%%%%%%%%%%%%%%%%

%----------------------------------------------------------------------------------------
%	PACKAGES AND OTHER DOCUMENT CONFIGURATIONS
%----------------------------------------------------------------------------------------

\documentclass[a0,portrait]{a0poster}

\usepackage{multicol} % This is so we can have multiple columns of text side-by-side
\columnsep=100pt % This is the amount of white space between the columns in the poster
\columnseprule=3pt % This is the thickness of the black line between the columns in the poster
\usepackage{fancybox}
\usepackage[svgnames]{xcolor} % Specify colors by their 'svgnames', for a full list of all colors available see here: http://www.latextemplates.com/svgnames-colors

\usepackage{epsfig}
\usepackage{pdfpages}
\usepackage{cite}
\usepackage{color}
\usepackage{colortbl}
\usepackage[cmex10]{amsmath}
\usepackage{amssymb}
\usepackage{soul}
\usepackage{amsfonts}
%\usepackage{mathtools}
\usepackage[normalem]{ulem}
%\graphicspath{{fig/}}
%\usepackage{soul}
\usepackage{booktabs}
\usepackage{color}
\usepackage{placeins}
\usepackage{upgreek}
\usepackage[utf8]{inputenc}

\usepackage[export]{adjustbox}

%\usepackage{times} % Use the times font
%\usepackage{palatino} % Uncomment to use the Palatino font
\usepackage{epstopdf}
\usepackage{graphicx} % Required for including images
%\graphicspath{{figures/}} % Location of the graphics files
\usepackage{booktabs} % Top and bottom rules for table
\usepackage[font=small,labelfont=bf]{caption} % Required for specifying captions to tables and figures
\usepackage{amsfonts, amsmath, amsthm, amssymb} % For math fonts, symbols and environments
\usepackage{wrapfig} % Allows wrapping text around tables and figures
\usepackage[framemethod=TikZ]{mdframed}
\usepackage{lipsum}
\usepackage{titlesec} % Modify titles
%\usepackage{fontspec}
%\setmainfont{Calibri}
\newcommand{\ds}{\displaystyle}
\newcommand {\bo}[1]{\textbf{#1}}
\newcommand{\pa}[1]{\left({#1}\right)}
\newcommand{\co}[1]{\left[{#1}\right]}
\newcommand{\ch}[1]{\left\{{#1}\right\}}

\titleformat{\section}{\color{white}\normalfont\Large\bfseries}{\color{white}\thesection}{1em}{\colorbox{SteelBlue}}{}

\setlength{\columnseprule}{0pt}
\mdfdefinestyle{MyFrame}{%
	linecolor=SteelBlue,
	outerlinewidth=2pt,
	roundcorner=50pt,
	innertopmargin=\baselineskip,
	innerbottommargin=\baselineskip,
	innerrightmargin=20pt,
	innerleftmargin=20pt,
	backgroundcolor=white!50!white}
%\input{figconfig}
%\titleformat{command}[shape]{format}{label}{sep}{be %fore-code}{after-code}
\begin{document}
\begin{mdframed}[style=MyFrame]
%----------------------------------------------------------------------------------------
%	POSTER HEADER 
%----------------------------------------------------------------------------------------

% The header is divided into two boxes:
% The first is 75% wide and houses the title, subtitle, names, university/organization and contact information
% The second is 25% wide and houses a logo for your university/organization or a photo of you
% The widths of these boxes can be easily edited to accommodate your content as you see fit
\begin{minipage}[b]{0.33\linewidth}
\raggedright
\includegraphics[width=4cm,valign=t]{Cinvestav-seal.pdf}
\end{minipage}
%
\begin{minipage}[b]{0.33\linewidth}
\centering
\includegraphics[width=5cm,height=5cm,valign=t]{neo_logo.png}
\end{minipage}
% 
\begin{minipage}[b]{0.33\linewidth}
\raggedleft
\includegraphics[width=5cm,valign=t]{Escudo-UNAM-escalable.pdf}
\end{minipage}\\

\vspace{3cm}
\begin{minipage}[h]{0.98\linewidth}
\centering \huge \color{SteelBlue} \textbf{On Objective Reduction of Many-objective Optimization by Means of PerformanceIndicators} \color{Black}\\ % Title
\Large \textbf{Fernando Avitúa Varela\textsuperscript{1}, Carlos Ignacio Hernández Castellanos\textsuperscript{2}}\\ % Author(s)
\normalsize \textsuperscript{1,2} Universidad Nacional Autónoma de México,\textsuperscript{1,2}\\ %[-0.5cm] % University/organization
fernando\_avitua@ciencias.unam.mx@email.com\textsuperscript{1}, carlos.hernandez@iimas.unam.mx\textsuperscript{2}\\
\end{minipage}
\vspace{0.5cm} % A bit of extra whitespace between the header and poster content

%----------------------------------------------------------------------------------------

\begin{multicols}{2} % This is how many columns your poster will be broken into, a portrait poster is generally split into 2 columns


%----------------------------------------------------------------------------------------
%	ABSTRACT
%----------------------------------------------------------------------------------------

\section{Abstract}
Indicator-Based Evolutionary Multi-Objective Algorithms (IB-EMOAs) have demonstrated efficacy in identifying Pareto Fronts where traditional Pareto dominance-based methods might not be as effective. However, selecting the appropriate indicator for specific problem contexts remains challenging. Decision-makers might prioritize a Pareto Front (PF) approximation that offers enhanced diversity rather than one with superior convergence. The PFI-EMOA algorithm has tackled this the trade-off between convergence and diversity by designing a composite indicator within the density estimator, leveraging a specific combination of the Inverted Generational Distance Plus (IGD+) and S-energy metrics. In our study, we explore the extent to which an IB-EMOA favors solutions that exhibit greater diversity or convergence. This is achieved by modifying the scalarization approach used in PFI-EMOA, systematically adjusting the weights assigned to IGD+ and S-energy. We then evaluate the impact of altering these weights on the overall quality of the Pareto Front approximation, as assessed by various Quality Indicators (QIs).

\section{Introduction}\label{section:1}

Quality Indicators (QIs) serve as critical tools for evaluating Pareto Front approximations, assessing metrics such as convergence to a reference set and the diversity of solutions [1]. Evolutionary Multi-Objective Algorithms (EMOAs) that rely on Pareto dominance often struggle with the identification of varied Pareto Fronts, especially as the number of objectives increases or the shape of the Pareto Front (PF) varies [2]. To overcome these limitations, Indicator-Based EMOAs (IB-EMOAs) have been successfully implemented. These algorithms employ QIs to navigate the search space, continuously balancing exploration and exploitation by evaluating the characteristics of the approximation set at each step of the MOEA process.

PFI-EMOA [3], a variant of IB-EMOA, parallels SMS-EMOA [4] but distinguishes itself through its density estimator, which incorporates a custom QI. This QI merges a convergence indicator (Inverted Generational Distance Plus, IGD+) and a diversity indicator (S-energy), employing an order-preserving Tchebyshev augmented function to integrate their contributions effectively. Initially, the algorithm employed fixed weight vectors w=(0.5,0.5) for scenarios with two and three objectives and w=(0.9,0.1) for configurations exceeding three objectives. This weighting strategy, favoring IGD+ for higher objective counts, was designed to mitigate the risk of convergence stagnation due to the progressively smaller dominated regions.

In this study, we conduct a statistical analysis to examine how varying these weight combinations—in favor of either convergence or diversity—affects the output of PFI-EMOA. We present our findings in Figure 1, highlighting the impact of different weighting strategies on the quality of the Pareto Front approximations.




\section{Methodology}\label{section2}
%
At vero eos et accusamus et iusto odio dignissimos ducimus qui blanditiis praesentium voluptatum deleniti atque corrupti quos dolores et quas molestias excepturi sint occaecati cupiditate non provident, similique sunt in culpa qui officia deserunt mollitia animi, id est laborum et dolorum fuga. Et harum quidem rerum facilis est et expedita distinctio. Nam libero tempore, cum soluta nobis est eligendi optio cumque nihil impedit quo minus id quod maxime placeat facere possimus, omnis voluptas assumenda est, omnis dolor repellendus. Temporibus autem quibusdam et aut officiis debitis aut rerum necessitatibus saepe eveniet ut et voluptates repudiandae sint et molestiae non recusandae. Itaque earum rerum hic tenetur a sapiente delectus, ut aut reiciendis voluptatibus maiores alias consequatur aut perferendis doloribus asperiores repellat.
%
\begin{equation}
    R = A B + \int_0^\infty x dx,
\end{equation}
%
onde $A$ e $B$ eos et accusamus et iusto odio dignissimos ducimus qui blanditiis praesentium voluptatum deleniti, respectivamente.

Et harum quidem rerum facilis est et expedita distinctio. Nam libero tempore, cum soluta nobis est eligendi optio cumque nihil impedit quo minus id quod maxime placeat facere possimus, omnis voluptas assumenda est, omnis dolor repellendus. Temporibus autem quibusdam et aut officiis debitis aut rerum necessitatibus saepe eveniet ut et voluptates repudiandae sint et molestiae non recusandae.
%
\begin{equation}\label{eqn:L}
\begin{split}
	f(x) = \frac{\sqrt{2 \pi}}{\int_{0}^{\infty} e^{-x^2} dx} \cdot \left(\sum_{n=1}^{10} \frac{\sin(nx)}{n^2} - \frac{\cos(2x)}{\log(x+1)}\right) + \frac{\Gamma\left(\frac{3}{4}\right)}{\lim_{{x \to 0}} \frac{e^x - 1}{x}},
\end{split}
\end{equation}
%
eta $\beta_i = {\pa{\eta^{2}-1}\mu\rho_{i}^{\alpha}}/{\left({2\eta}\right)}$.

Sed ut perspiciatis unde omnis iste natus error sit voluptatem accusantium doloremque laudantium, totam rem aperiam, eaque ipsa quae ab illo inventore veritatis et quasi architecto beatae vitae dicta sunt explicabo. Nemo enim ipsam voluptatem quia voluptas sit aspernatur aut odit aut fugit, sed quia consequuntur magni dolores eos qui ratione voluptatem sequi nesciunt. Neque porro quisquam est, qui dolorem ipsum quia dolor sit amet, consectetur, adipisci velit, sed quia non numquam eius modi tempora incidunt ut labore et dolore magnam aliquam quaerat voluptatem. Ut enim ad minima veniam, quis nostrum exercitationem ullam corporis suscipit laboriosam, nisi ut aliquid ex ea commodi consequatur? Quis autem vel eum iure reprehenderit qui in ea voluptate velit esse quam nihil molestiae consequatur, vel illum qui dolorem eum fugiat quo voluptas nulla pariatur?

\begin{align}
    F(x) &= \frac{\sqrt{e^{2x}}}{\sin^2(x) + \cos^2(x)} \cdot \left(\sum_{n=0}^{+\infty} \frac{(2x)^n}{n!} - \int_{0}^{\frac{\pi}{2}} \frac{e^{2x}}{x^2 + \arctan^2(x)} \, dx\right) \nonumber \\
        &\times\frac{\log_{10}\left(\frac{e^{3x}}{\sqrt[3]{x+1}}\right)}{\sqrt[4]{\frac{2x}{x^2 + 1}}} \cdot \frac{\Gamma\left(\frac{3}{4}\right) + \frac{1}{2} \cdot \ln\left(\frac{x^2 + 1}{x+1}\right)}{\lim_{{x \to 0}} \frac{e^x - 1}{x}}.
\end{align}

Lorem ipsum dolor sit amet, consectetur adipiscing elit, sed do eiusmod tempor incididunt ut labore et dolore magna aliqua. Ut enim ad minim veniam, quis nostrud exercitation ullamco laboris nisi ut aliquip ex ea commodo consequat. Duis aute irure dolor in reprehenderit in voluptate velit esse cillum dolore eu fugiat nulla pariatur. Excepteur sint occaecat cupidatat non proident, sunt in culpa qui officia deserunt mollit anim id est laborum.

At vero eos et accusamus et iusto odio dignissimos ducimus qui blanditiis praesentium voluptatum deleniti atque corrupti quos dolores et quas molestias excepturi sint occaecati cupiditate non provident, similique sunt in culpa qui officia deserunt mollitia animi, id est laborum et dolorum fuga.
%
\begin{equation}\label{eqn:L1}
\begin{split}
	f(x) = \frac{\sqrt{2 \pi}}{\int_{0}^{\infty} e^{-x^2} dx} \cdot \left(\sum_{n=1}^{10} \frac{\sin(nx)}{n^2} - \frac{\cos(2x)}{\log(x+1)}\right) + \frac{\Gamma\left(\frac{3}{4}\right)}{\lim_{{x \to 0}} \frac{e^x - 1}{x}},
\end{split}
\end{equation}

\section{Results}\label{section4}

At vero eos et accusamus et iusto odio dignissimos ducimus qui blanditiis praesentium voluptatum deleniti atque corrupti quos dolores et quas molestias excepturi sint occaecati cupiditate non provident, similique sunt in culpa qui officia deserunt mollitia animi, id est laborum et dolorum fuga.
%
\vspace{1cm}
%\hspace{1cm}
\begin{minipage}{\columnwidth}
    \includegraphics[width=0.9\columnwidth]{f}
\end{minipage}
%\hspace{1cm}
\vspace{1cm} %\\\vspace{-10cm}

%\hspace{1cm}
\begin{minipage}{\columnwidth}
    \includegraphics[width=0.9\columnwidth]{g}
\end{minipage}
%\hspace{1cm}

\section{Conclusions and future work}

Sed ut perspiciatis unde omnis iste natus error sit voluptatem accusantium doloremque laudantium, totam rem aperiam, eaque ipsa quae ab illo inventore veritatis et quasi architecto beatae vitae dicta sunt explicabo. Nemo enim ipsam voluptatem quia voluptas sit aspernatur aut odit aut fugit, sed quia consequuntur magni dolores eos qui ratione voluptatem sequi nesciunt. Neque porro quisquam est, qui dolorem ipsum quia dolor sit amet, consectetur, adipisci velit, sed quia non numquam eius modi tempora incidunt ut labore et dolore magnam aliquam quaerat voluptatem. Ut enim ad minima veniam, quis nostrum exercitationem ullam corporis suscipit laboriosam, nisi ut aliquid ex ea commodi consequatur? Quis autem vel eum iure reprehenderit qui in ea voluptate velit esse quam nihil molestiae consequatur, vel illum qui dolorem eum fugiat quo voluptas nulla pariatur?

Lorem ipsum dolor sit amet, consectetur adipiscing elit, sed do eiusmod tempor incididunt ut labore et dolore magna aliqua. Ut enim ad minim veniam, quis nostrud exercitation ullamco laboris nisi ut aliquip ex ea commodo consequat. Duis aute irure dolor in reprehenderit in voluptate velit esse cillum dolore eu fugiat nulla pariatur. Excepteur sint occaecat cupidatat non proident, sunt in culpa qui officia deserunt mollit anim id est laborum.

At vero eos et accusamus et iusto odio dignissimos ducimus qui blanditiis praesentium voluptatum deleniti atque corrupti quos dolores et quas molestias excepturi sint occaecati cupiditate non provident, similique sunt in culpa qui officia deserunt mollitia animi, id est laborum et dolorum fuga.

\section*{Agradecimentos}
%
Este trabalho foi parcialmente suportado por XXXXXX, com recursos  YYYYY.
%----------------------------------------------------------------------------------------
%	REFERENCES
%----------------------------------------------------------------------------------------

%\bibliographystyle{IEEEtran}
%\bibliography{IEEEabrv,references}

%---------------------------------------------------
-------------------------------------
\end{multicols}
\begin{center}
\color{SteelBlue}{SBrT 2024 - Bel\'{e}m, PA - 01 a 04 de Outubro de 2024}
\end{center}
\end{mdframed}
\end{document}